\begin{abstract}
Contextual advertising buys on the page instead of the person, so a bidder must know, inside each bid, what the reader is looking at. On real ad-exchange traffic that means millions of distinct pages a day, a large share of requests on pages never seen before, and slots whose surroundings differ from the page as a whole. We present \sys{}, a system that labels pages against the 704 categories of the IAB Content Taxonomy and serves the labels to a bidder in microseconds. Offline, a crawler ordered by traffic feeds a budgeted cascade of five tiers, from the URL alone through a lexicon, text embeddings and image embeddings to a small language model. A page's budget follows its traffic, each tier runs only while the page is in doubt, and each label keeps the quote and reason behind it. Online, a library reads the profile, backs off to the page's section and site, and, when a page module reports the blocks around a slot, computes the slot's own context: relevance blends the page with nearby content through a kernel set per ad format, while suitability takes the worst case over everything on screen with the ad. On a golden set of 383 pages in four main languages, the tuned cascade reaches a hierarchical F1 of 0.501 while calling the language model on 29.8\% of pages, at about \$42 per million pages. On bid logs, buyers bid about twice as often on pages that match their deal's context. The golden labels were judged by a language model, and the slot model's weights are priors. We give a pre-registered experiment, with sample sizes, that would test page against slot context per format and fit how far context reaches. The code, a made-up golden set in the same format, and this paper are at \repo{} (MIT licence). The real golden set comes from an exchange's traffic and is not published.
\end{abstract}
